\textbf{Dépôt :} \url{https://github.com/DARK-SHAD0W/Clusterisation-Conteneurs}

\begin{itemize}
  \item \code{mfp/} : l'application My Favorite Places (client React, API Node.js/Express, base PostgreSQL) et le cluster Docker Swarm (\code{mfp/swarm/}) ;
  \item \code{voting-app/} : l'application de vote du TD 3.
\end{itemize}

\textbf{Environnement de travail :}

\begin{center}
\small
\begin{tabular}{l l p{6.2cm}}
  \toprule
  Outil & Version & Remarque \\
  \midrule
  Docker Engine & 29.7.2 & dans WSL 2 (Ubuntu 24.04), sans Docker Desktop \\
  Docker Compose & v5.5.0 & plugin \code{docker compose} \\
  Docker des nœuds du cluster & 29.8.2 & image \code{docker:dind} \\
  Node.js / Yarn & 22.23.2 / 1.22.22 & dans WSL \\
  Bruno & 4.2.1 & client API, sous Windows \\
  \bottomrule
\end{tabular}
\end{center}

Les commandes du cluster se lancent depuis le dossier \code{mfp/swarm/}, qui contient son \code{compose.yml}.
